% SPDX-License-Identifier: CC-BY-SA-4.0
% Author: Matthieu Perrin
% Part: 
% Section: 
% Sub-section: 
% Frame: 

\begingroup

\begin{frame}{Conditions de progression}

  \begin{description}[\alert{Progrès non-bloquant :}]
  \item[\color{exampleColor}Progrès global :] les rapides peuvent toujours passer devant les lents
  \item[\color{exampleColor}Progrès local :] équité sur la terminaison 
  \item[\alert{Progrès bloquant :}] les threads lents peuvent gêner les rapides
  \item[\alert{Progrès non-bloquant :}] tolérance aux pannes
  \end{description}

  \pause
  
  \begin{center}
    \begin{tabular}{|r|c|c|}
      \hline
      & \example{Progrès global} & \example{Progrès local}\\
      \hline
      \alert{Progrès bloquant}  & Deadlock-free & Starvation-free\\
      \hline
      \alert{Progrès non-bloquant} & Lock-free     & Wait-free\\
      \hline
    \end{tabular}
  \end{center}

  \begin{block}{Définitions}
    Toutes les opérations des threads corrects se terminent \structure{sous l'hypothèse} :
    \begin{description}[Starvation-freedom :]
    \item[Deadlock-freedom :] \alert{il n'y a pas de crash} et \example{le nombre d'appels est fini}
    \item[Starvation-freedom :] \alert{il n'y a pas de crash}
    \item[Lock-freedom :] \example{le nombre d'appels est fini}
    \item[Wait-freedom :] vrai
    \end{description}
  \end{block}

  \footnoteref{Maurice Herlihy, Nir Shavit. \emph{On the nature of progress.} PODS (2011)}
\end{frame}

\endgroup
\endinput
